% ===================================================================
%  GAMBAR No. 12 — Setasiun. --- Sepur. --- Kapal.
%
%  Traduction, faite en 2026, en javanais d'aujourd'hui, sur l'ido de
%  texto/io. Regles de la colonne : voir 10-gambar-01.tex. Une seule
%  numerotation.
%
%  LE SEUL « \textsc{} » DU LIVRET HORS DIALOGUE EST DANS CE TABLEAU,
%  ET C'EST LUI QUI A FIXE LE SEUIL DE kolonoj.py. L'ido y compose
%  « \textsc{Noto.} --- » pour l'avertissement sur les heures : une
%  attribution de forme, mais pas de locuteur. Le tableau 5 en compte
%  trente-six, celui-ci UNE, les quatorze autres aucune ; le seuil de
%  cinq laisse donc les deux cas de part et d'autre sans rien serrer,
%  et les regles de niveau restent actives ici comme il faut.
%
%  LE CHEMIN DE FER EST LE TABLEAU LE PLUS NEERLANDAIS DU LIVRET, ET
%  DE LOIN. sepur la voie et le train, saka la gare — de statie —,
%  karcis le wagon, lokomotip, rel, wesel l'aiguillage, sinyal, peron,
%  kondhektur, masinis, bagasi, kelas, restorasi, tunnel. Le javanais
%  n'a pas un seul mot a lui pour cette machine, et c'est juste :
%  Java a eu son premier chemin de fer en 1867, construit par la
%  Nederlandsch-Indische Spoorweg Maatschappij, et le mot « sepur »
%  est « spoor ».
%
%  ET « sepur » EST DEVENU PLUS JAVANAIS QUE LE MOT INDONESIEN. La
%  langue nationale dit « kereta api », voiture-a-feu, un compose
%  malais transparent ; le javanais dit « sepur », un mot neerlandais
%  qu'aucun Javanais ne sent plus etranger, et qui a donne des verbes
%  et des composes — nyepur, sepuran. C'est le quatrieme couple apres
%  potlot/pensil, pit/sepeda et porok/garpu, et le seul ou la forme
%  javanaise soit un emprunt CONTRE une forme indonesienne native.
%  L'emprunt n'est pas un signe de faiblesse : il est un signe de date.
%
%  LE TELEGRAPHE (5) ET L'HORLOGE (89) SONT LES DEUX SEULES CHOSES DU
%  TABLEAU QUE LE JAVANAIS NOMME SANS EMPRUNTER LE MOT ENTIER : jam
%  gedhe, jarum gedhe (90), jarum cilik (91) — le meme jam et les
%  memes jarum qu'au clocher du tableau 3. Le temps se disait ici
%  avant le train.
%
%  ET LE VOYAGE EST EN ALLEMAGNE, AVEC UNE FRONTIERE D'AVANT 1914 QUE
%  L'IDO SIGNALE LUI-MEME EN NOTE. La colonne la garde telle quelle :
%  on modernise la langue, jamais les choses — et une frontiere est
%  une chose.
% ===================================================================

%%K t12-apar-1 apar
\VUsaut{4.0mm}
\VUtitre{15.0pt}{60}{GAMBAR No. 12}
\VUsaut{2.4mm}
\VUpk{11.4pt}{40}{Setasiun. --- Sepur. --- Kapal.}
\VUsaut{3.4mm}
\textsc{Cathetan.} --- \textit{Tabel jam kang dianggo ing saben negara
beda-beda ; mula kita nganggo tuladha jam ing} \VUgras{tabel Prancis.}

%%K t12-01-1 p
1. --- Aku lagi wae lelungan liwat Jerman. Sadurunge mangkat, aku tuku
\VUgras{buku jam} \textsuperscript{(15)} supaya ngerti jam mangkate sepur.
Sawise mesthekake yen sepur kang paling trep bakal mangkat saka Paris
jam sanga bengi lan tekan Strasburg jam siji luwih telulas menit, aku
nyiyapake lelunganku, lan ing dina esuke aku dieterake menyang setasiun.

%%K t12-02-1 p
2. --- Nalika aku mlebu aula mangkat kang gedhe, ing mburine
\VUgras{pandhita} \textsuperscript{(2)} desa tuwa, kang nyangking
\VUgras{tas lelungan}e \textsuperscript{(3)}, akeh \VUgras{wong kang arep
lelungan} \textsuperscript{(1)} wis padha teka.

%%K t12-02-2 p
Sarehne aku kepengin ngirim \VUgras{telegram} \textsuperscript{(5)}, aku
njaluk \VUgras{juru priksa} \textsuperscript{(6)} nuduhake \VUgras{kantor
telegraf} \textsuperscript{(4)} marang aku.

%%K t12-03-1 p
3. --- Sadurunge mlebu \VUgras{ruwang ngenteni} \textsuperscript{(7)}, aku
lunga tuku \VUgras{karcis} \textsuperscript{(9)} pulang-balik.

%%K t12-04-1 p
4. --- Aku ora suwe ngenteni ing \VUgras{loket} \textsuperscript{(8)},
amarga mung sethithik wong ing kono. \VUgras{Prajurit}
\textsuperscript{(10)}, kang arep mangkat prei, kanthi grapyak menehake
gilirane marang aku, mula aku duwe wektu cukup kanggo njaluk sawetara
katrangan marang \VUgras{kepala setasiun} \textsuperscript{(13)}, kang
lagi ngobrol karo \VUgras{komisaris} \textsuperscript{(12)} pengawasan lan
karo \VUgras{polisi} \textsuperscript{(11)} kang lagi tugas.

%%K t12-tit-1 sub
\VUsaut{3.0mm}
\VUpk{10.2pt}{40}{Ndhaptarake bagasi.}
\VUsaut{3.0mm}

%%K t12-05-1 p
5. --- Sawise tuku koran ing \VUgras{kios buku} \textsuperscript{(14)},
aku njaluk \VUgras{bagasi}ku \textsuperscript{(17)} ditimbang, kang lagi
wae digawa \VUgras{kuli} \textsuperscript{(16)} nganggo \VUgras{gerobag
barang} \textsuperscript{(19)} ; \VUgras{juru tulis} \textsuperscript{(22)}
kang ngurusi \VUgras{timbangan} \textsuperscript{(21)} awèh ngerti marang
aku yen kopherku aboté telung puluh lima kilogram — dadi ana
\VUgras{luwih bobot} limang kilogram, kang kudu dakbayar tambahan ing
\VUgras{loket} bagasi \textsuperscript{(23)}.

%%K t12-06-1 p
6. --- Sarehne aku kesuwen tekaku, aku duwe wektu longgar kanggo
nyawang mlebu-metune para wong lelungan ing setasiun. Sawetara
nyelehake utawa njupuk maneh bagasi ciliké ing \VUgras{papan nitipake
bagasi} \textsuperscript{(25)} ; liyane maca \VUgras{tabel jam}
\textsuperscript{(26)}. Wong lelungan kang teka gage-gage menyang
\VUgras{lawang metu} \textsuperscript{(32)} karo \VUgras{koper}e
\textsuperscript{(24)} ing tangan. Sawetara ngenteni ing \VUgras{papan
mbagi bagasi} \textsuperscript{(27)} lan nuduhake \VUgras{karcis}e
\textsuperscript{(29)} kanggo njupuk maneh kopher, \VUgras{peti}
\textsuperscript{(28)} utawa \VUgras{kranjang}e \textsuperscript{(30)}.

%%K t12-07-1 p
7. --- \VUgras{Para petugas cukai} \textsuperscript{(33)} kanthi
ngati-ati mriksa bungkusan supaya mesthekake apa ana barang kang kudu
dilaporake.

%%K t12-08-1 p
8. --- Sawetara wong lelungan menyang \VUgras{restorasi}
\textsuperscript{(34)}, panggonan \VUgras{pelayan} \textsuperscript{(35)}
kanthi sregep ngladeni. Wong-wong iku kudu gage-gage, amarga
\VUgras{sepur} \textsuperscript{(36)} iku sepur ekspres, ora suwe mandheg
ing setasiun.

%%K t12-09-1 p
9. --- Banjur aku menyang peron mangkat lan nggoleki \VUgras{gerbong}
\textsuperscript{(37)} langsung supaya ora kepeksa ganti gerbong ing
tengah dalan. Aku munggah menyang kreta koridor kang duwe
\VUgras{ruwangan} \textsuperscript{(38)} kanggo wong udud lan ruwangan
kang dicawisake kanggo « ibu-ibu piyambakan ».

%%K t12-tit-2 sub
\VUsaut{3.0mm}
\VUpk{10.2pt}{40}{Mangkat ing wayah bengi.}
\VUsaut{3.0mm}

%%K t12-10-1 p
10. --- Sawise munggah ing \VUgras{undhakan} \textsuperscript{(40)},
nutup \VUgras{lawang} \textsuperscript{(39)}, nggulung \VUgras{tirai}
\textsuperscript{(41)} lan nyelehake bagasi ciliku ing \VUgras{rak
jaring} \textsuperscript{(43)}, aku lungguh ing \VUgras{bangku}
\textsuperscript{(42)}. Nalika weruh \VUgras{tukang lampu}
\textsuperscript{(44)} ngurubake lampu, aku mikir yen aku bakal
nglampahi sewengi ing gerbong, banjur aku ngundang petugas kang ngurusi
nyewakake \VUgras{bantal} \textsuperscript{(45)} kanggo njaluk siji.

%%K t12-11-1 p
11. --- Kondhektur sepur teka mriksa karcis. Aku takon marang dheweke
kena apa kita durung mangkat, amarga jame wis pas. Dheweke mangsuli yen
\VUgras{kepala sepur} \textsuperscript{(46)} lagi repot nyelehake pit ing
\VUgras{gerbong barang} \textsuperscript{(47)}. Kejaba iku, kabeh
pemanas sikil \textsuperscript{(48)} durung diganti.

%%K t12-12-1 p
12. --- Nanging kita sedhela maneh bakal mangkat, amarga \VUgras{tukang
geni} \textsuperscript{(50)}, ing \VUgras{tender}e \textsuperscript{(49)},
wis njupuk areng kanggo nguripake genine \VUgras{lokomotip}
\textsuperscript{(54)}, lan \VUgras{masinis} \textsuperscript{(51)} mung
ngenteni tandhane \VUgras{wakil kepala setasiun} \textsuperscript{(52)}
kang ana ing \VUgras{peron} \textsuperscript{(53)}.

%%K t12-13-1 p
13. --- \VUgras{Ketel} \textsuperscript{(56)} lokomotip iku nggereng
alon ; \VUgras{cerobong}e \textsuperscript{(55)} ngutahake kukus kang
kecampur \VUgras{uwab} \textsuperscript{(60)}. \VUgras{Sempritan}
\textsuperscript{(59)} kang nyoceki muni seru, lan \VUgras{piston}
\textsuperscript{(58)} wiwit obah, nyurung \VUgras{rodha}
\textsuperscript{(57)} mesin kang abot iku.

%%K t12-tit-3 sub
\VUsaut{3.0mm}
\VUpk{10.2pt}{40}{Mangkat!}
\VUsaut{3.0mm}

%%K t12-14-1 p
14. --- Rikate sepur, kang mlaku ing \VUgras{rel} \textsuperscript{(64)},
saya banter ; \VUgras{peron} \textsuperscript{(65)} ilang ; kita ngliwati
\VUgras{kakus} \textsuperscript{(66)}, lan uga rangkeyan gerbong kang
diparkir ing \VUgras{sepur} sijine \textsuperscript{(63)} :
\VUgras{gerbong turu} \textsuperscript{(67)}, \VUgras{gerbong restorasi}
\textsuperscript{(68)}, \VUgras{gerbong pos} \textsuperscript{(69)}.

%%K t12-noto-1 noto
\textit{(*) Cathetan.} --- \textit{Mesthine gegambaran lelungan iki
miturut tapel wates sadurunge perang.}

%%K t12-15-1 p
15. --- Banjur, \VUgras{setasiun barang} \textsuperscript{(71)} liwat ing
ngarepe mripat kita. Ing sangisore hangar, para petugas ngemot
\VUgras{barang} \textsuperscript{(72)} kang dikirim nganggo sepur alon.
\VUgras{Wong regu langsir} \textsuperscript{(76)} ing cedhake
\VUgras{tandhon banyu} \textsuperscript{(73)} nglangsir \VUgras{gerbong
kewan} \textsuperscript{(75)} lan \VUgras{gerbong lantai rata}
\textsuperscript{(79)} kanggo nyusun \VUgras{sepur barang}
\textsuperscript{(80)}.

%%K t12-16-1 p
16. --- Kita ngliwati \VUgras{persimpangan} \textsuperscript{(78)} kang
pungkasan, panggonan tukang wesel ngadeg ; kita ngliwati \VUgras{piring
sinyal} \textsuperscript{(86)}, lan setasiune ilang ing kadohan.

%%K t12-17-1 p
17. --- Sedhela maneh kita ngliwati \VUgras{kreteg wesi}
\textsuperscript{(74)} kang nyabrangi \VUgras{kali gedhe}
\textsuperscript{(81)}. Sawise ngliwati kreteg iku, kita ngetutake kali,
panggonan \VUgras{kapal panarik} \textsuperscript{(82)} kang kuwat narik
\VUgras{prau momotan} \textsuperscript{(83)} kang abot. Kita uga weruh
\VUgras{kapal panumpang} \textsuperscript{(84)} nalika lagi ngrapet ing
\VUgras{ponton} \textsuperscript{(85)}.

%%K t12-18-1 p
18. --- Sawise ngliwati sawetara setasiun kanthi banter banget, kita
ngliwati sawetara \VUgras{trowongan} \textsuperscript{(87)} lan ninggal
ing mburi \VUgras{omah cilik} kang tanpa wilangan \textsuperscript{(88)}
duweke \VUgras{penjaga palang.} Kabuyung dening goyange sepur, aku turu
kepati.

%%K t12-tit-4 sub
\VUsaut{3.0mm}
\VUpk{10.2pt}{40}{Setasiun Deutsch-Avricourt.}
\VUsaut{3.0mm}

%%K t12-19-1 p
19. --- Aku dumadakan tangi dening swarane petugas kang mbengok :
« Deutsch-Avricourt! \textsuperscript{(*)}. Kabeh mudhun! » Banjur ana
pamriksan pabean lan kabeh urusan tapel wates kang njelehi. Aku kudu
nyocogake jam kanthongku karo \VUgras{jam gedhe} \textsuperscript{(89)}
setasiun, kang luwih dhisik sèket lima menit ; lagi wae tangi, aku
angel mbedakake \VUgras{jarum gedhe} \textsuperscript{(90)} karo
\VUgras{kang cilik} \textsuperscript{(91)} ; \VUgras{papan angka}ne
\textsuperscript{(92)} dhewe blas ora cetha merga pedhut esuk.

%%K t12-20-1 p
20. --- Nalika para petugas pabean nindakake pamriksan, aku karo angop
maca \VUgras{plakat} \textsuperscript{(93)} kang ngrengga temboke
setasiun. Aku sarapan kanggo nglong wektu, ndeleng peta \VUgras{jaringan}
\textsuperscript{(94)} sepur Jerman kanggo ngerti sepira adohe kang isih
misahake aku karo Strasburg, banjur aku mlebu maneh menyang gerbongku,
kakuwatake dening pangarep-arep sedhela maneh tekan tujuane lelunganku.

\VUsaut{6.0mm}
\VUfilet{22mm}
